\begin{figure}[htbp]
\centering
\begin{tikzpicture}[x=1cm,y=1cm,>=Latex,font=\small]
\begin{scope}[xshift=1.8cm,yshift=.3cm]
\draw[->,gray] (0,-.5)--(0,3.65) node[above] {\(n/9\)};
\foreach \k in {0,1,2,3}{
 \draw[gray!60,dashed] (-1.1,\k)--(1.1,\k);
 \node[left,font=\scriptsize] at (-1.1,\k){\(\k\)};
}
\draw[black,line width=.65pt,domain=0:27,samples=260,smooth,variable=\n]
 plot ({cos(40*\n)},{\n/9+.14*sin(40*\n)});
\foreach \n in {0,...,27}{
 \fill ({cos(40*\n)},{\n/9+.14*sin(40*\n)}) circle (.026);
}
\foreach \k in {0,1,2,3}{\fill (1,\k) circle (.045);}
\node[align=center,text width=3.3cm,font=\scriptsize] at (0,-1.02)
 {Same phase at \(n=0,9,18,27\);\\different memory degree.};
\end{scope}
\begin{scope}[xshift=4.2cm,yshift=.5cm]
\node[anchor=west] at (0,3.1){\(z_n(0)=\lfloor(n+1)\delta\rfloor-\lfloor n\delta\rfloor\)};
\draw[->] (0,1.7)--(9,1.7) node[right]{\(n\)};
\foreach \n in {21,43,65,87,109,131,152,174,196,218}{
 \draw[line width=.6pt] ({\n*.039},1.7)--({\n*.039},2.25);
 \fill ({\n*.039},2.25) circle (.033);
 \node[below,font=\tiny] at ({\n*.039},1.7){\n};
}
\foreach \a/\b/\g in {21/43/22,43/65/22,65/87/22,87/109/22,109/131/22,131/152/21,152/174/22,174/196/22,196/218/22}{
 \node[font=\scriptsize] at ({(\a+\b)*.0195},2.55){\g};
}
\node[anchor=west,align=left,text width=8.3cm] at (0,.55)
 {Irrational coding determines the intervals between events. Return of the nonadic phase does not turn this sequence into a periodic word.};
\end{scope}
\end{tikzpicture}
\caption{\raggedright Complementary representations of transport. Left: projection of the helix \((\cos(2\pi n/9),\sin(2\pi n/9),n/9)\): advancing nine steps recovers the phase and increases the height. Right: the first ten events of the lower word for \(\delta=\log_{10}(10/9)\), with exact separations \(21\) or \(22\). The helical representation introduces no additional physical metric.}
\label{fig:holonomia-esturmiana}
\end{figure}
